\section{Executive Summary}
\subsection{Overview}
% <EDIT> One paragraph describing what was assessed and the headline outcome.
A white-box static security assessment was conducted against \reportsystem. \fillme{One sentence
on the sweep methodology, e.g. "the assessment swept every file for the full OWASP-style
vulnerability catalogue, and each Critical/High/Medium finding was adversarially re-verified."}

The assessment identified \textbf{\fillme{N} distinct findings}. \fillme{One or two sentences
summarising the overall posture and the worst realistic outcome.}

\begin{criticalbox}
\textbf{Result --- \fillme{one-line verdict, e.g. "the module is critically insecure"}.}\\[6pt]
Headline conclusions:
\begin{itemize}
  \item \fillme{Headline finding 1}
  \item \fillme{Headline finding 2}
  \item \fillme{Headline finding 3}
  \item \fillme{N of M findings require no authentication.}
\end{itemize}
\end{criticalbox}

\subsection{Overall Risk Rating}
\begin{center}\begin{tikzpicture}
  \node[fill=AccentRed, text=white, font=\Large\bfseries\sffamily, minimum width=9cm, minimum height=1.2cm, rounded corners=4pt]{CRITICAL};
\end{tikzpicture}\end{center}

\subsection{Findings by Severity}
\begin{center}\begin{tabularx}{\linewidth}{@{}l r X@{}}
\toprule\rowcolor{TableHead}
\textcolor{white}{\bfseries Severity} & \textcolor{white}{\bfseries Count} & \textcolor{white}{\bfseries Meaning} \\
\midrule
\critical & \fillme{n} & \fillme{e.g. Unauthenticated RCE / SQLi / auth bypass} \\
\rowcolor{TableAlt}
\high & \fillme{n} & \fillme{e.g. Authenticated RCE/SQLi, stored XSS, IDOR exposing PII} \\
\medium & \fillme{n} & \fillme{e.g. Reflected XSS, CSRF, missing session hardening} \\
\rowcolor{TableAlt}
\low & \fillme{n} & \fillme{e.g. Hygiene / defence-in-depth gaps} \\
\midrule
\multicolumn{1}{@{}l}{\textbf{Total}} & \textbf{\fillme{N}} & \fillme{n unauthenticated} \\
\bottomrule\end{tabularx}\end{center}

\subsection{Findings by Vulnerability Class}
\begin{center}\begin{tabularx}{\linewidth}{@{}l r X@{}}
\toprule\rowcolor{TableHead}
\textcolor{white}{\bfseries Class} & \textcolor{white}{\bfseries Count} & \textcolor{white}{\bfseries Severities} \\
\midrule
\fillme{e.g. SQL Injection} & \fillme{n} & \fillme{n Critical, n Medium} \\
\rowcolor{TableAlt}
\fillme{e.g. Broken Access Control / IDOR} & \fillme{n} & \fillme{n High} \\
\fillme{e.g. Sensitive Data Exposure} & \fillme{n} & \fillme{n High, n Medium, n Low} \\
\bottomrule\end{tabularx}\end{center}

\subsection{Remediation Roadmap (priorities)}
\begin{center}\begin{tabularx}{\linewidth}{@{}l X@{}}
\toprule\rowcolor{TableHead}
\textcolor{white}{\bfseries Priority} & \textcolor{white}{\bfseries Action} \\
\midrule
\textbf{P1 --- Stop the worst outcome} & \fillme{e.g. Remediate any injection sinks reaching the OS/filesystem.} \\
\rowcolor{TableAlt}
\textbf{P2 --- Systemic} & \fillme{e.g. Parameterise all SQL; add a central auth filter; output-encode for XSS; add CSRF tokens.} \\
\textbf{P3 --- Process} & \fillme{e.g. Externalise + encrypt secrets; rotate exposed credentials; hash passwords correctly.} \\
\bottomrule\end{tabularx}\end{center}
